% Luc Destombe --- licence CC BY-NC-SA 4.0
% https://creativecommons.org/licenses/by-nc-sa/4.0/deed.fr
% Fichier autonome : compiler avec pdflatex, deux fois (signets, nombre de pages).
\def\lyceeversion{1.0}
\documentclass[11pt,a4paper]{article}
\makeatletter
\newif\ifeleve
\elevefalse

\newif\iflycee@cm
\newif\iflycee@palatino

\RequirePackage[utf8]{inputenc}
\RequirePackage[T1]{fontenc}
\RequirePackage[french]{babel}
\RequirePackage{etoolbox}

\newcommand{\ShorthandsOff}{\@ifundefined{shorthandoff}{}{\shorthandoff{;:!?}}}
\newcommand{\ShorthandsOn}{\@ifundefined{shorthandon}{}{\shorthandon{;:!?}}}
\ShorthandsOff
\AtBeginDocument{\ShorthandsOn}
\AtBeginEnvironment{tikzpicture}{\ShorthandsOff}

\RequirePackage{geometry}
\RequirePackage{fancyhdr}
\RequirePackage{lastpage}
\RequirePackage{multicol}
\RequirePackage{array}
\RequirePackage{enumitem}
\RequirePackage{xcolor}
\RequirePackage{needspace}

\RequirePackage{amsmath,amssymb}
\iflycee@cm\else\RequirePackage{fourier}\fi
\DeclareMathAlphabet{\mathcal}{OMS}{cmsy}{m}{n}
\renewcommand{\ttdefault}{lmtt}
\begingroup\catcode`\_=\active \catcode`\^=\active
\global\let_\sb \global\let^\sp
\endgroup
\AtBeginDocument{\catcode`\_=12 \mathcode`\_="8000
  \catcode`\^=12 \mathcode`\^="8000 }
\newlength{\retraitcalculs}
\setlength{\retraitcalculs}{2em}

\RequirePackage{tikz}
\usetikzlibrary{arrows.meta,calc,babel,shapes.misc}
\RequirePackage[most]{tcolorbox}
\tcbuselibrary{skins,breakable}

\definecolor{coulDef}{HTML}{1F4E79}
\definecolor{coulProp}{HTML}{7030A0}
\definecolor{coulRem}{HTML}{C55A11}
\definecolor{coulEx}{HTML}{2E7D32}
\definecolor{coulExo}{HTML}{B03A2E}
\definecolor{coulPart}{HTML}{1F4E79}
\definecolor{coulMeth}{HTML}{1B6E4F}
\definecolor{coulCourbe}{HTML}{0B5ED7}
\definecolor{coulCourbeB}{HTML}{C00000}
\definecolor{coulCourbeC}{HTML}{2E7D32}
\definecolor{coulGrille}{HTML}{8C8C8C}
\definecolor{coulRep}{HTML}{1B6E4F}
\definecolor{coulBar}{HTML}{2E7D32}

\newcommand{\R}{\mathbb{R}}
\newcommand{\Rs}{\mathbb{R}^{*}}
\newcommand{\N}{\mathbb{N}}
\newcommand{\Ns}{\mathbb{N}^{*}}
\newcommand{\Z}{\mathbb{Z}}
\newcommand{\Q}{\mathbb{Q}}
\renewcommand{\le}{\leqslant}
\renewcommand{\ge}{\geqslant}
\renewcommand{\approx}{\simeq}
\newcommand{\vect}[1]{\overrightarrow{#1}}
\newcommand{\norme}[1]{\left\lVert #1 \right\rVert}
\newcommand{\intervalle}[2]{\left[#1\,;#2\right]}
\RequirePackage{cancel}
\renewcommand{\CancelColor}{\color{coulExo}}
\newcommand{\fois}[1]{\times{\color{coulExo}#1}}
\newcommand{\barre}[1]{\cancel{#1}}
\newlength{\lycee@eqh}\newlength{\lycee@eqd}
\newsavebox{\lycee@eqa}\newsavebox{\lycee@eqb}\newsavebox{\lycee@eqc}
\newcommand{\lycee@eqcol}[1]{\hbox{$\renewcommand{\arraystretch}{1.35}%
  \begin{array}[t]{@{}>{\displaystyle}r@{}>{\displaystyle{}}l@{}}#1\end{array}$}}
\newcommand{\lycee@eqmesure}[1]{%
  \ifdim\ht#1>\lycee@eqh\setlength{\lycee@eqh}{\ht#1}\fi
  \ifdim\dp#1>\lycee@eqd\setlength{\lycee@eqd}{\dp#1}\fi}
\newcommand{\lycee@eqtrait}{\hspace{1.2em}\tikz[baseline=0pt]\draw[dash pattern=on 2pt off 2pt,
  line width=0.5pt] (0,\lycee@eqh) -- (0,-\lycee@eqd);\hspace{1.2em}}
\newcommand{\eqproduit}[2]{\par\smallskip\noindent C'est une équation produit :
  $A\times B=0 \iff A=0$ ou $B=0$.\par\nopagebreak\smallskip
  \begingroup\setlength{\lycee@eqh}{0pt}\setlength{\lycee@eqd}{0pt}%
  \setbox\lycee@eqa=\lycee@eqcol{#1}\setbox\lycee@eqb=\lycee@eqcol{#2}%
  \lycee@eqmesure\lycee@eqa\lycee@eqmesure\lycee@eqb%
  \noindent\hspace*{\retraitcalculs}\usebox\lycee@eqa\lycee@eqtrait\usebox\lycee@eqb\endgroup\par\smallskip}
\newcommand{\eqproduittrois}[3]{\par\smallskip\noindent C'est une équation produit :
  $A\times B\times C=0 \iff A=0$ ou $B=0$ ou $C=0$.\par\nopagebreak\smallskip
  \begingroup\setlength{\lycee@eqh}{0pt}\setlength{\lycee@eqd}{0pt}%
  \setbox\lycee@eqa=\lycee@eqcol{#1}\setbox\lycee@eqb=\lycee@eqcol{#2}\setbox\lycee@eqc=\lycee@eqcol{#3}%
  \lycee@eqmesure\lycee@eqa\lycee@eqmesure\lycee@eqb\lycee@eqmesure\lycee@eqc%
  \noindent\hspace*{\retraitcalculs}\usebox\lycee@eqa\lycee@eqtrait\usebox\lycee@eqb\lycee@eqtrait\usebox\lycee@eqc\endgroup
  \par\smallskip}

\tikzset{grille/.style={coulGrille,line width=0.4pt}}
\newcommand{\quadrillage}[4]{\draw[grille,step=1] (#1,#2) grid (#3,#4);}
\tikzset{
  croix/.style={cross out,draw,line width=0.6pt,inner sep=0pt,minimum size=4pt},
  croixdroite/.style={croix,rotate=45,line width=0.8pt},
}
\newcommand{\pt}[3]{\path (#1) node[croix]{} node[#2,font=\small,inner sep=2.5pt]{$#3$};}

\newcommand{\lycee@repere}[4]{%
  \draw[grille,step=1] (#1,#2) grid (#3,#4);
  \draw[-{Stealth[length=2mm]},thick] (#1,0)--({#3+0.4},0) node[below left=-1pt]{$x$};
  \draw[-{Stealth[length=2mm]},thick] (0,#2)--(0,{#4+0.4}) node[below left=-1pt]{$y$};
  \node[below left,font=\scriptsize,inner sep=1.5pt] at (0,0) {$0$};}
\newcommand{\lycee@gradx}[1]{\foreach \k/\l in {#1}{%
  \draw (\k,0.07)--(\k,-0.07) node[below,font=\scriptsize,inner sep=1.5pt]{$\l$};}}
\newcommand{\lycee@grady}[1]{\foreach \k/\l in {#1}{%
  \draw (0.07,\k)--(-0.07,\k) node[left,font=\scriptsize,inner sep=1.5pt]{$\l$};}}
\newcommand{\lycee@intff}[2]{\left[#1\,;#2\right]}
\newcommand{\lycee@intfo}[2]{\left[#1\,;#2\right[}
\newcommand{\lycee@intof}[2]{\left]#1\,;#2\right]}
\newcommand{\lycee@intoo}[2]{\left]#1\,;#2\right[}
\newcommand{\lycee@ens}[1]{\left\{#1\right\}}
\AtBeginDocument{%
  \providecommand{\repere}{\lycee@repere}%
  \providecommand{\gradx}{\lycee@gradx}%
  \providecommand{\grady}{\lycee@grady}%
  \providecommand{\Cf}{\mathcal{C}\sb{f}}%
  \providecommand{\Cg}{\mathcal{C}\sb{g}}%
  \providecommand{\intff}{\lycee@intff}%
  \providecommand{\intfo}{\lycee@intfo}%
  \providecommand{\intof}{\lycee@intof}%
  \providecommand{\intoo}{\lycee@intoo}%
  \providecommand{\ens}{\lycee@ens}}

\newlist{questions}{enumerate}{3}
\setlist[questions,1]{label=\arabic*\textdegree),leftmargin=*,itemsep=2pt,topsep=3pt}
\setlist[questions,2]{label=\alph*),leftmargin=*,itemsep=1pt,topsep=2pt}
\setlist[questions,3]{label=\roman*),leftmargin=*,itemsep=1pt,topsep=1pt}
\setlist[itemize]{label=\textbullet,leftmargin=*,itemsep=2pt,topsep=3pt}

\newcommand{\besoinplace}[1]{\par\penalty\@M\begingroup
  \dimen@=#1\relax\dimen@ii=\pagegoal\advance\dimen@ii-\pagetotal
  \ifdim\dimen@>\dimen@ii\ifdim\dimen@ii>\z@\vfil\fi\break\fi\endgroup}

\newcommand{\lycee@pied}{\thepage/\pageref{LastPage}}
\newcommand{\entete}[2]{%
  \pagestyle{fancy}\fancyhf{}%
  \fancyhead[L]{\footnotesize\color{coulPart}\textbf{#1}}%
  \fancyhead[R]{\footnotesize\color{coulPart}\textbf{#2}}%
  \fancyfoot[C]{\footnotesize\lycee@pied}%
  \renewcommand{\headrulewidth}{0.4pt}%
  \renewcommand{\headrule}{\hbox to\headwidth{%
    \color{coulPart!40}\leaders\hrule height \headrulewidth\hfill}}}

\RequirePackage{ccicons}
\newcommand{\lycee@licence}{{\scriptsize\color{gray}%
  \copyright\ Luc Destombe --- \ccbyncsa\ CC BY-NC-SA 4.0}}
\newif\iflycee@licence \lycee@licencetrue
\newcommand{\sanslicence}{\lycee@licencefalse}
\AtBeginDocument{\iflycee@licence\AddToHookNext{shipout/foreground}{%
  \put(0,0){\rlap{\kern\dimexpr1in+\hoffset+\oddsidemargin+\textwidth\relax
    \llap{\raisebox{-\dimexpr1in+\voffset+\topmargin+\headheight+\headsep
      +\textheight+\footskip\relax}{\lycee@licence}}}}}\fi}

\newcommand{\lycee@site}{faire-des-maths.fr}
\newcommand{\lycee@version}{{\scriptsize\color{gray}\lycee@site
  \ifdefined\lyceeversion\ --- Version~\lyceeversion\fi}}
\AtBeginDocument{\iflycee@licence\AddToHookNext{shipout/foreground}{%
  \put(0,0){\rlap{\kern\dimexpr1in+\hoffset+\oddsidemargin\relax
    \raisebox{-\dimexpr1in+\voffset+\topmargin+\headheight+\headsep
      +\textheight+\footskip\relax}{\lycee@version}}}}\fi}

\newcommand{\formulesserrees}{%
  \setlength{\abovedisplayskip}{3pt}\setlength{\belowdisplayskip}{3pt}%
  \setlength{\abovedisplayshortskip}{2pt}\setlength{\belowdisplayshortskip}{3pt}}

\setlength{\parindent}{0pt}

\RequirePackage[implicit=false,hidelinks,pdfencoding=auto,psdextra,bookmarksopen,
  bookmarksopenlevel=1,pdfcreator={},pdfproducer={}]{hyperref}
\RequirePackage{bookmark}
\hypersetup{pdfauthor={Luc Destombe},pdfkeywords={CC BY-NC-SA 4.0}}
\newcounter{lycee@signet}
\newcommand{\signet}[3][]{\stepcounter{lycee@signet}%
  \edef\lycee@dest{\if\relax\detokenize{#1}\relax
    signet.\arabic{lycee@signet}\else#1\fi}%
  \hypertarget{\lycee@dest}{}\bookmark[level=#2,dest=\lycee@dest]{#3}}
\pdfstringdefDisableCommands{%
  \def\R{R}\def\N{N}\def\Z{Z}\def\Q{Q}\def\vect#1{#1}\def\Cf{Cf}\def\Cg{Cg}\def\fois#1{×#1}%
  \def\barre#1{#1}\def\intervalle#1#2{[#1;#2]}\def\intff#1#2{[#1;#2]}%
  \def\textsuperscript#1{#1}\def\dfrac#1#2{#1/#2}\def\frac#1#2{#1/#2}%
  \def\vec#1{#1⃗}}%
\newcommand{\res}[1]{#1}
\geometry{top=2cm,bottom=1cm,left=1cm,right=1cm,headheight=14pt,footskip=0.6cm}

\RequirePackage{pgfplots}
\pgfplotsset{compat=1.18}
\RequirePackage{tkz-tab}

\newcounter{partiecours}
\renewcommand{\thepartiecours}{\Roman{partiecours}}
\newcommand{\partiecours}[1]{%
  \refstepcounter{partiecours}%
  \Needspace*{8\baselineskip}%
  \bigskip
  \noindent\begin{tcolorbox}[enhanced,colback=coulPart,colframe=coulPart,
      boxrule=0pt,arc=2pt,left=8pt,right=8pt,top=4pt,bottom=4pt,
      before skip=6pt,after skip=10pt]
    \color{white}\bfseries\large \thepartiecours{} --- #1%
    \signet[partie-\arabic{partiecours}]{1}{\thepartiecours{} --- #1}
  \end{tcolorbox}%
  \addcontentsline{toc}{section}{\thepartiecours{} --- #1}%
}

\newtcolorbox{boitecours}[3][]{%
  enhanced,breakable,
  colback=white,colframe=#2,coltitle=white,
  fonttitle=\bfseries,
  attach boxed title to top left={xshift=8pt,yshift=-\tcboxedtitleheight/2},
  boxed title style={colback=#2,arc=2pt,boxrule=0pt},
  boxrule=0.9pt,arc=2pt,
  left=8pt,right=8pt,top=8pt,bottom=6pt,
  title={#3},#1}

\newcounter{methode}
\newenvironment{definitionbox}[1][Définition]%
  {\begin{boitecours}[phantom={\signet{2}{#1}}]{coulDef}{#1}}{\end{boitecours}}
\newenvironment{proprietebox}[1][Propriété]%
  {\begin{boitecours}[phantom={\signet{2}{#1}}]{coulProp}{#1}}{\end{boitecours}}
\newenvironment{methodebox}[1][Méthode]{\stepcounter{methode}%
  \begin{boitecours}[phantom={\signet[methode-\arabic{methode}]{2}{#1}}]{coulMeth}{#1}}%
  {\end{boitecours}}

\newcommand{\etape}[1]{\par\smallskip\textbf{\color{coulMeth}#1}\par\nobreak\smallskip}
\newcommand{\enonce}[1]{\emph{#1}\par\smallskip}
\newcommand{\reponse}[1]{\par\smallskip
  {\footnotesize\itshape\color{black!55}Réponses : #1}}
\newcommand{\demo}[1]{\textbf{\color{coulMeth}Démonstration.} #1}
\newcommand{\afaire}[1]{%
  \par\smallskip
  \noindent{\small\color{coulExo}$\blacktriangleright$\ \textbf{À faire :}
    fiche d'exercices du chapitre, #1.}\par\smallskip}

\newenvironment{calculs}{\setlength{\jot}{5pt}\@fleqntrue\@mathmargin\retraitcalculs
  \setlength{\abovedisplayskip}{4pt}\setlength{\belowdisplayskip}{4pt}%
  \setlength{\abovedisplayshortskip}{4pt}\setlength{\belowdisplayshortskip}{4pt}%
  \csname alignat*\endcsname{2}}%
  {\csname endalignat*\endcsname}
\newcommand{\rem}[1]{\text{\small\itshape\color{black!60}#1}}
\newcommand{\explique}[2]{\par\smallskip\noindent
  \begin{minipage}[t]{0.57\linewidth}\raggedright #1\end{minipage}\hfill
  \begin{minipage}[t]{0.40\linewidth}\raggedright\leavevmode
    {\small\itshape\color{black!60}#2}\end{minipage}\par}
\newcommand{\cas}[1]{\par\smallskip\noindent\textbf{\color{coulExo}#1}\vadjust{\nobreak}}
\newcommand{\calc}[1]{\par\smallskip\textbf{\color{coulExo}#1}\par\nobreak}

\newtcolorbox{remarquebox}[1][Remarque]{%
  enhanced,breakable,colback=white,colframe=coulRem,
  boxrule=0pt,leftrule=2.5pt,arc=0pt,outer arc=0pt,
  left=8pt,right=6pt,top=5pt,bottom=5pt,
  before upper={\textbf{\color{coulRem}#1}\par\smallskip}}

\newtcolorbox{exemplebox}[1][Exemple]{%
  enhanced,breakable,colback=white,colframe=coulEx,
  boxrule=0pt,leftrule=2.5pt,arc=0pt,outer arc=0pt,
  left=8pt,right=6pt,top=5pt,bottom=5pt,
  before upper={\textbf{\color{coulEx}#1}\par\smallskip}}

\pgfplotsset{
  repere/.style={
    axis lines=middle,
    axis line style={-{Stealth[length=2.4mm]},thick},
    every axis x label/.style={at={(ticklabel* cs:1.0)},anchor=west},
    every axis y label/.style={at={(ticklabel* cs:1.0)},anchor=south},
    label style={font=\footnotesize},
    tick label style={font=\scriptsize},
    grid=both,
    major grid style={grille},
    minor tick num=0,
    tick align=inside,
    clip mode=individual,
  },
}
\tikzset{
  courbe/.style={coulCourbe,line width=1.1pt,smooth},
  courbeB/.style={coulCourbeB,line width=1.1pt,smooth},
  courbeC/.style={coulCourbeC,line width=1.1pt,smooth},
}

\newlist{etapes}{enumerate}{1}
\setlist[etapes]{label=\textbf{\arabic*.},leftmargin=*,itemsep=1pt,topsep=2pt}

\newlist{expressions}{enumerate}{1}
\setlist[expressions]{label={\Alph*\ $=$},leftmargin=*,itemsep=3pt,topsep=3pt}

\newcounter{exo}
\renewcommand{\theexo}{\ifnum\value{exo}<10 0\fi\arabic{exo}}
\newtcolorbox{exercice}[1][]{%
  enhanced,breakable,
  colback=white,colframe=coulExo,
  boxrule=0.9pt,arc=2pt,
  left=8pt,right=8pt,top=8pt,bottom=6pt,
  coltitle=white,fonttitle=\bfseries,
  attach boxed title to top left={xshift=8pt,yshift=-\tcboxedtitleheight/2},
  boxed title style={colback=coulExo,arc=2pt,boxrule=0pt},
  phantom={\signet{2}{Exercice \arabic{exo}}},
  title={Exercice \theexo\ifx\relax#1\relax\else{} --- #1\fi}}
\newenvironment{exo}[1][\relax]{\refstepcounter{exo}\begin{exercice}[#1]}{\end{exercice}}

  \newenvironment{correction}{\par}{\par}
  \newcommand{\autableau}{}
  \newcommand{\etapeexemples}{\etape{Exemples corrigés}}
  \newcommand{\etapetableau}[1]{\etape{#1}}
  \newenvironment{exempletableau}[1][Exemple]%
    {\begin{exemplebox}[#1]}{\end{exemplebox}}

\RequirePackage{comment}
\excludecomment{correctionavj}

\newcommand{\coursEntete}{}
\newcommand{\coursNiveau}{}
\newcommand{\coursTitre}{}
\newcommand{\coursSurTitre}{}
\newcommand{\coursSousTitre}{}

\newcommand{\enteteCours}[1]{\renewcommand{\coursEntete}{#1}}
\newcommand{\chapitrecours}[2]{\enteteCours{Chapitre #1 --- #2}}
\newcommand{\niveaucours}[1]{\renewcommand{\coursNiveau}{#1}}
\newcommand{\titrecours}[3][]{%
  \renewcommand{\coursTitre}{#2}%
  \renewcommand{\coursSurTitre}{#1}%
  \renewcommand{\coursSousTitre}{#3}}

\pagestyle{fancy}
\fancyhf{}
\fancyhead[L]{\footnotesize\color{coulPart}\textbf{\coursEntete}}
\fancyhead[R]{\footnotesize\color{coulPart}\textbf{\coursNiveau}}
\fancyfoot[C]{\footnotesize\thepage}
\renewcommand{\headrulewidth}{0.4pt}
\renewcommand{\headrule}{\hbox to\headwidth{\color{coulPart!40}\leaders\hrule height \headrulewidth\hfill}}

\setlength{\parskip}{2pt}

\newcommand{\lycee@mentioneleve}{\\[2pt]{\small\itshape Version élève}}
\AtBeginDocument{%
  \begin{center}
    \begin{tcolorbox}[enhanced,width=0.72\linewidth,colback=white,colframe=coulPart,
        boxrule=1.6pt,arc=3pt,halign=center,top=10pt,bottom=10pt]
      \ifx\coursSurTitre\empty
        {\Huge\bfseries\color{coulPart} \coursTitre}\\[4pt]
      \else
        {\Huge\bfseries\color{coulPart} \coursTitre}\\[3pt]
        {\Large\bfseries\color{coulPart} \coursSurTitre}\\[5pt]
      \fi
      {\normalsize \coursSousTitre}%
      
    \end{tcolorbox}
  \end{center}%
}

\RequirePackage{tikz}
\RequirePackage{xcolor}

\usetikzlibrary{decorations.text}
\definecolor{lycee@metal}{HTML}{1F4E79}
\definecolor{lycee@manche}{HTML}{8A5A2B}
\definecolor{lycee@etincelle}{HTML}{F28C28}
\newcommand{\lycee@pluriel}[2]{#1 #2\ifnum#1>1 s\fi}
\newcommand{\lycee@sceau}{%
  \edef\lycee@bas{\lycee@pluriel{\lyceerelectures}{relecture} et
    \lycee@pluriel{\lyceerodages}{rodage}}%
  \begin{tikzpicture}[line join=round]
    \fill[coulPart!10,draw=coulPart,line width=.8pt] (0:2.2)
      \foreach \i in {1,...,71}{ -- (\i*5:{isodd(\i)?2.1:2.2}) } -- cycle;
    \draw[coulPart,line width=1.2pt] (0,0) circle (1.93);
    \draw[coulPart,line width=.4pt] (0,0) circle (1.85);
    \draw[coulPart,line width=.4pt] (0,0) circle (1.2);
    \path[decorate,decoration={text along path,text align=center,
        text={|\fontsize{9.5}{11}\selectfont\bfseries\color{coulPart}|TEST{\'E} AVEC MES {\'E}L{\`E}VES}}]
      (180:1.45) arc (180:0:1.45);
    \path[decorate,decoration={text along path,text align=center,
        text/.expanded={|\noexpand\fontsize{8.5}{10}\noexpand\selectfont\noexpand\color{coulPart}|\lycee@bas}}]
      (195:1.68) arc (195:345:1.68);
    \begin{scope}[x=.25mm,y=-.25mm]
      \begin{scope}[shift={(-32,-40)}]%
        \fill[lycee@metal] (23,31) rectangle (41,49);
        \draw[lycee@etincelle,line width=.5pt,line cap=round]
          (20,30)--(15,27) (19,34)--(14,35) (44,30)--(49,27) (45,34)--(50,35);
        \fill[lycee@manche,rounded corners=.3mm] (38,24.6) rectangle (60,27.6);
        \fill[lycee@metal,rounded corners=.4mm] (25,20) rectangle (39,31);
        \fill[lycee@etincelle,rounded corners=.2mm] (25,20) rectangle (28,31);
      \end{scope}
    \end{scope}
    \ifdefined\lyceesceauplusieurs\def\lycee@annee{Ann{\'e}es~:}\else
      \def\lycee@annee{Ann{\'e}e}\fi
    \path[decorate,decoration={text along path,text align=center,
        text/.expanded={|\noexpand\fontsize{8}{10}\noexpand\selectfont\noexpand\color{coulPart}|%
          \unexpanded\expandafter{\lycee@annee} \unexpanded\expandafter{\lyceesceauannees}}}]
      (95:1.08) arc (95:445:1.08);
  \end{tikzpicture}}
\AtEndDocument{\ifdefined\lyceesceauannees
  \par\besoinplace{5cm}\vspace{1.5em}%
  \begin{center}\lycee@sceau\end{center}\fi}
\makeatother

\chapitrecours{2}{Fonctions et pourcentages}
\niveaucours{1\textsuperscript{re} STMG}
\titrecours{FONCTIONS ET POURCENTAGES}{Cours --- Première STMG}

\begin{document}

\begin{remarquebox}[Comment utiliser ce cours]
Chaque partie commence par ce qu'il faut \textbf{savoir} (définitions, propriétés), puis
vient ce qu'il faut \textbf{savoir faire} : une méthode, avec des
\textbf{exemples corrigés} faits en classe, puis des exercices
\og\textbf{À vous de jouer}\fg{} à chercher seul. Tout le chapitre se fait
\textbf{sans calculatrice}.

\end{remarquebox}

\partiecours{Tableau de valeurs, courbe et lecture graphique}

\begin{remarquebox}[Image et antécédent (rappel de seconde)]
Si $f(x)=y$, on dit que $y$ est l'\textbf{image} de $x$ et que $x$ est un
\textbf{antécédent} de $y$. Un nombre a une seule image ; il peut avoir plusieurs
antécédents, ou aucun.
\end{remarquebox}

\begin{definitionbox}[Définition 1 --- Courbe représentative]
La \textbf{courbe représentative} $\Cf$ de $f$ est formée des points $M\bigl(x\,;f(x)\bigr)$.
Autrement dit, le point $M(x\,;y)$ est sur $\Cf$ si et seulement si $y=f(x)$.
\end{definitionbox}

\begin{methodebox}[Méthode 1 --- Tracer une courbe à partir d'un tableau de valeurs]
\etapeexemples
Tracer la courbe de la fonction $f$ définie sur $\intff{-2}{4}$ par $f(x)=x^{2}-2x-3$.
\begin{correction}
\cas{Résolution}\par\nopagebreak
\begin{minipage}[c]{0.58\linewidth}\raggedright
\explique{On calcule des images :}{un tableau de valeurs}
\begin{center}\small
\renewcommand{\arraystretch}{1.25}
\begin{tabular}{|c|c|c|c|c|c|c|c|}
\hline
$x$ & $-2$ & $-1$ & $0$ & $1$ & $2$ & $3$ & $4$\\
\hline
$f(x)$ & $5$ & $0$ & $-3$ & $-4$ & $-3$ & $0$ & $5$\\
\hline
\end{tabular}
\end{center}
\explique{On place les points $\bigl(x\,;f(x)\bigr)$, puis on les relie par une courbe
régulière, sans règle.}{par exemple $(1\,;-4)$, le point le plus bas}
\end{minipage}\hfill
\begin{minipage}[c]{0.38\linewidth}\centering
\begin{tikzpicture}[x=0.38cm,y=0.38cm]
  \repere{-3}{-5}{5}{5}
  \gradx{-2,-1,1,2,3,4} \grady{-4,-3,-2,-1,1,2,3,4}
  \draw[coulCourbe,line width=1.1pt,domain=-2:4,samples=60] plot (\x,{\x*\x-2*\x-3});
  \foreach \p in {(-2,5),(-1,0),(0,-3),(1,-4),(2,-3),(3,0),(4,5)}
    \path[coulCourbe] \p node[croix]{};
  \node[coulCourbe,font=\footnotesize] at (4.3,3) {$\Cf$};
\end{tikzpicture}
\end{minipage}
\end{correction}

\etape{À vous de jouer}
Tracer la courbe de la fonction $g$ définie sur $\intff{-2}{4}$ par $g(x)=-x^{2}+2x+3$.
\reponse{$g(-2)=-5$ ;\; $g(-1)=0$ ;\; $g(0)=3$ ;\; $g(1)=4$ ;\; $g(2)=3$ ;\;
$g(3)=0$ ;\; $g(4)=-5$ ;\; point le plus haut $(1\,;4)$.}
\end{methodebox}

\begin{methodebox}[Méthode 2 --- Lire une image, lire des antécédents]
\etapeexemples
Sur la courbe de la méthode 1 :
\begin{questions}[label=\alph*)]
  \item lire l'image de $1$ ;
  \item lire les antécédents de $0$.
\end{questions}
\begin{correction}
\begin{minipage}[c]{0.58\linewidth}\raggedright
\cas{a)}
\explique{$f(1)=\res{-4}$.}{de $1$ sur l'axe des abscisses, on va verticalement jusqu'à la
courbe, puis on lit l'ordonnée (en vert)}
\cas{b)}
\explique{Les antécédents de $0$ sont \res{$-1$ et $3$}.}{on trace l'horizontale $y=0$ et on
lit les abscisses de \textbf{tous} les points de la courbe qui sont dessus (en rouge)}
\end{minipage}\hfill
\begin{minipage}[c]{0.38\linewidth}\centering
\begin{tikzpicture}[x=0.38cm,y=0.38cm]
  \repere{-3}{-5}{5}{5}
  \gradx{-2,-1,1,2,3,4} \grady{-4,-3,-2,-1,1,2,3,4}
  \draw[coulCourbe,line width=1.1pt,domain=-2:4,samples=60] plot (\x,{\x*\x-2*\x-3});
  \draw[coulEx,dashed,thick] (1,0) -- (1,-4) -- (0,-4);
  \path[coulEx] (1,-4) node[croix]{};
  \path[coulExo] (-1,0) node[croix]{} (3,0) node[croix]{};
  \node[coulCourbe,font=\footnotesize] at (4.3,3) {$\Cf$};
\end{tikzpicture}
\end{minipage}
\end{correction}

\etape{À vous de jouer}
Sur la même courbe, lire :
\begin{questions}
  \item l'image de $-2$, puis celle de $2$ ;
  \item les antécédents de $-3$ ;
  \item les antécédents de $-4$, puis ceux de $-5$.
\end{questions}
\reponse{1\textdegree) $f(-2)=5$ et $f(2)=-3$ ;\; 2\textdegree) $0$ et $2$ ;\;
3\textdegree) $-4$ a un seul antécédent, $1$ ; aucun pour $-5$, la courbe ne descend
pas jusque-là.}
\end{methodebox}

\afaire{exercices 1 à 3}

\partiecours{Tableau de variations d'une fonction}

\begin{definitionbox}[Définition 2 --- Fonction croissante, décroissante]
Sur un intervalle $I$ :
\begin{itemize}
  \item $f$ est \textbf{croissante} si, quand $x$ augmente, $f(x)$ augmente : la courbe
        monte ; l'ordre est conservé ;
  \item $f$ est \textbf{décroissante} si, quand $x$ augmente, $f(x)$ diminue : la courbe
        descend ; l'ordre est renversé.
\end{itemize}
\begin{center}
\begin{minipage}{0.46\linewidth}\centering
{\small\textbf{$f$ croissante}}\par\smallskip
\begin{tikzpicture}[x=0.6cm,y=0.6cm]
  \draw[grille,step=1] (0,0) grid (4,4);
  \draw[-{Stealth[length=2mm]},thick] (0,0) -- (4.5,0) node[below left=-1pt] {$x$};
  \draw[-{Stealth[length=2mm]},thick] (0,0) -- (0,4.5) node[below left=-1pt] {$y$};
  \draw[coulCourbe,line width=1.1pt,domain=0.2:4,samples=40] plot (\x,{0.4+0.22*\x*\x});
  \draw[coulEx,dashed] (1,0) -- (1,0.62) -- (0,0.62);
  \draw[coulEx,dashed] (3,0) -- (3,2.38) -- (0,2.38);
  \path[coulEx] (1,0.62) node[croix]{} (3,2.38) node[croix]{};
  \node[below,font=\scriptsize] at (1,0) {$x_{1}$};
  \node[below,font=\scriptsize] at (3,0) {$x_{2}$};
  \node[left,font=\scriptsize] at (0,0.62) {$f(x_{1})$};
  \node[left,font=\scriptsize] at (0,2.38) {$f(x_{2})$};
\end{tikzpicture}\par
{\footnotesize $x_{1}<x_{2}$ donne $f(x_{1})<f(x_{2})$}
\end{minipage}\hfill
\begin{minipage}{0.46\linewidth}\centering
{\small\textbf{$f$ décroissante}}\par\smallskip
\begin{tikzpicture}[x=0.6cm,y=0.6cm]
  \draw[grille,step=1] (0,0) grid (4,4);
  \draw[-{Stealth[length=2mm]},thick] (0,0) -- (4.5,0) node[below left=-1pt] {$x$};
  \draw[-{Stealth[length=2mm]},thick] (0,0) -- (0,4.5) node[below left=-1pt] {$y$};
  \draw[coulCourbe,line width=1.1pt,domain=0.2:4,samples=40] plot (\x,{3.9-0.22*\x*\x});
  \draw[coulEx,dashed] (1,0) -- (1,3.68) -- (0,3.68);
  \draw[coulEx,dashed] (3,0) -- (3,1.92) -- (0,1.92);
  \path[coulEx] (1,3.68) node[croix]{} (3,1.92) node[croix]{};
  \node[below,font=\scriptsize] at (1,0) {$x_{1}$};
  \node[below,font=\scriptsize] at (3,0) {$x_{2}$};
  \node[left,font=\scriptsize] at (0,3.68) {$f(x_{1})$};
  \node[left,font=\scriptsize] at (0,1.92) {$f(x_{2})$};
\end{tikzpicture}\par
{\footnotesize $x_{1}<x_{2}$ donne $f(x_{1})>f(x_{2})$}
\end{minipage}
\end{center}
\end{definitionbox}

\begin{definitionbox}[Définition 3 --- Tableau de variations, extremums]
Le \textbf{tableau de variations} résume le sens de variation : une flèche qui monte
quand $f$ est croissante, une flèche qui descend quand elle est décroissante, et les
valeurs de $f$ au bout de chaque flèche.

\smallskip
Sur un intervalle, le \textbf{maximum} de $f$ est la plus grande valeur de $f(x)$, le
\textbf{minimum} la plus petite. On dit toujours \textbf{en quel $x$} il est atteint.
\end{definitionbox}

\begin{methodebox}[Méthode 3 --- Dresser un tableau de variations à partir d'une courbe]
\etapeexemples
La fonction $f$ est définie sur $\intff{-4}{4}$ par sa courbe ci-dessous. Dresser son
tableau de variations, puis donner son maximum et son minimum.
\begin{center}
\begin{tikzpicture}[x=0.42cm,y=0.42cm]
  \repere{-5}{-4}{5}{4}
  \gradx{-4,-3,-2,-1,1,2,3,4} \grady{-3,-2,-1,1,2,3}
  \draw[coulCourbe,line width=1.1pt]
    (-4,-3) .. controls (-3.4,-0.6) and (-2.6,3) .. (-2,3)
            .. controls (-1,3) and (1,-2) .. (2,-2)
            .. controls (2.8,-2) and (3.6,0.4) .. (4,2);
  \foreach \p in {(-4,-3),(-2,3),(2,-2),(4,2)} \path[coulCourbe] \p node[croix]{};
  \node[coulCourbe,font=\footnotesize] at (4.4,3) {$\Cf$};
\end{tikzpicture}
\end{center}
\begin{correction}
\cas{Résolution}\par\nopagebreak
\explique{La courbe monte de $x=-4$ à $x=-2$, descend jusqu'à $x=2$, puis remonte
jusqu'à $x=4$.}{on repère où la courbe change de sens}
\explique{$f(-4)=-3$ ; \ $f(-2)=3$ ; \ $f(2)=-2$ ; \ $f(4)=2$.}{on lit les valeurs aux
bornes et aux changements de sens}
\begin{center}
\begin{tikzpicture}
  \tkzTabInit[lgt=1.6,espcl=2.4]{$x$/1,$f$/2}{$-4$,$-2$,$2$,$4$}
  \tkzTabVar{-/$-3$,+/$3$,-/$-2$,+/$2$}
\end{tikzpicture}
\end{center}
\explique{Maximum : \res{$3$, atteint en $x=-2$}. \\ Minimum : \res{$-3$, atteint en $x=-4$}.}{on
compare toutes les valeurs du bas des flèches ; le minimum peut être à une borne}
\end{correction}

\etape{À vous de jouer}
La fonction $g$ est définie sur $\intff{-5}{3}$ par sa courbe ci-dessous.
\begin{center}
\begin{tikzpicture}[x=0.42cm,y=0.42cm]
  \repere{-6}{-3}{4}{5}
  \gradx{-5,-4,-3,-2,-1,1,2,3} \grady{-2,-1,1,2,3,4}
  \draw[coulCourbe,line width=1.1pt]
    (-5,4) .. controls (-4,1.5) and (-2.2,-2) .. (-1,-2)
           .. controls (0.2,-2) and (2.2,0.8) .. (3,3);
  \foreach \p in {(-5,4),(-1,-2),(3,3)} \path[coulCourbe] \p node[croix]{};
  \node[coulCourbe,font=\footnotesize] at (3.3,4) {$\Cg$};
\end{tikzpicture}
\end{center}
\begin{questions}
  \item Dresser le tableau de variations de $g$.
  \item Donner le maximum et le minimum de $g$.
\end{questions}
\reponse{1\textdegree) flèche qui descend de $4$ à $-2$, puis qui monte jusqu'à $3$ ;\;
2\textdegree) maximum $4$ en $x=-5$ ; minimum $-2$ en $x=-1$.}
\end{methodebox}

\afaire{exercices 4 à 7}

\partiecours{Appliquer un pourcentage}

\begin{definitionbox}[Définition 4 --- Pourcentage]
$t\,\%$ est une autre façon d'écrire le nombre $\dfrac{t}{100}$. Par exemple :
\[
  35\,\%=\dfrac{35}{100}=0{,}35 \qquad\quad 6\,\%=\dfrac{6}{100}=0{,}06 \qquad\quad
  120\,\%=\dfrac{120}{100}=1{,}2
\]
Prendre $t\,\%$ d'une quantité, c'est la \textbf{multiplier} par $\dfrac{t}{100}$.
\end{definitionbox}

\begin{remarquebox}[Pour calculer de tête]
\begin{center}
\renewcommand{\arraystretch}{1.4}
\begin{tabular}{|c|c|c|c|c|c|}
\hline
$1\,\%$ & $10\,\%$ & $5\,\%$ & $25\,\%$ & $50\,\%$ & $75\,\%$\\
\hline
diviser par $100$ & diviser par $10$ & moitié de $10\,\%$ & le quart &
la moitié & trois quarts\\
\hline
\end{tabular}
\end{center}
\end{remarquebox}

\begin{methodebox}[Méthode 4 --- Calculer $t\,\%$ d'une quantité]
\etapeexemples
\begin{questions}[label=\alph*)]
  \item Calculer $40\,\%$ de $120$.
  \item Une salle de concert a $1\,200$ places ; $15\,\%$ sont réservées aux abonnés.
        Combien de places sont réservées ?
  \item Dans une ville, $30$ vélos en libre-service représentent $5\,\%$ du parc.
        Combien le parc compte-t-il de vélos ?
\end{questions}
\begin{correction}
\cas{a)}
\explique{$10\,\%$ de $120$ font $12$.}{on divise par $10$}
\explique{$40\,\%$ de $120$ font $4\times 12=\res{48}$.}{$40\,\%$, c'est $4$ fois
$10\,\%$}
\cas{b)}
\explique{$10\,\%$ de $1\,200$ font $120$ ; $5\,\%$ en font $60$.}{$5\,\%$ : la moitié
de $10\,\%$}
\explique{$15\,\%$ de $1\,200$ font $120+60=180$ : \res{$180$ places} sont réservées.}{}
\cas{c)}
\explique{$5\,\%$ est la vingtième partie de $100\,\%$.}{$20\times 5=100$}
\explique{Le parc compte $20\times 30=\res{600}$ vélos.}{}
\end{correction}

\etape{À vous de jouer}
\begin{questions}
  \item Calculer $25\,\%$ de $360$ ; $10\,\%$ de $48$ ; $75\,\%$ de $60$.
  \item Dans une classe de $30$ élèves, $40\,\%$ font du sport en club. Combien
        d'élèves cela fait-il ?
  \item $18$ personnes représentent $25\,\%$ d'un groupe. Combien le groupe
        compte-t-il de personnes ?
\end{questions}
\reponse{$90$ ;\; $4{,}8$ ;\; $45$ ;\; $12$ élèves ;\; $72$ personnes.}
\end{methodebox}

\begin{methodebox}[Méthode 5 --- Augmenter ou diminuer une quantité de $t\,\%$]
\etapeexemples
\begin{questions}[label=\alph*)]
  \item Une console à $300$~€ augmente de $10\,\%$. Calculer son nouveau prix.
  \item Un manteau à $120$~€ est soldé à $-25\,\%$. Calculer son prix soldé.
\end{questions}
\begin{correction}
\cas{a)}
\explique{$10\,\%$ de $300$ font $30$.}{le montant de la hausse}
\explique{Nouveau prix : $300+30=\res{330}$~€.}{une hausse s'ajoute}
\cas{b)}
\explique{$25\,\%$ de $120$ font $30$.}{le quart de $120$ : le montant de la baisse}
\explique{Prix soldé : $120-30=\res{90}$~€.}{une baisse se retire}
\end{correction}

\etape{À vous de jouer}
\begin{questions}
  \item Un abonnement de $400$~€ augmente de $3\,\%$. Calculer le nouveau prix.
  \item Une commune de $8\,000$ habitants perd $5\,\%$ de sa population. Combien
        d'habitants compte-t-elle ensuite ?
  \item Un sac à $60$~€ est soldé à $-30\,\%$. Calculer son prix soldé.
\end{questions}
\reponse{$412$~€ ;\; $7\,600$ habitants ;\; $42$~€.}
\end{methodebox}

\afaire{exercices 8 à 12}

\partiecours{Calculer un pourcentage}

\begin{definitionbox}[Définition 5 --- Proportion]
Dans un ensemble de $N$ individus (le « total »), une partie en compte $n$. La
\textbf{proportion} de cette partie est
\[
  \boxed{p=\dfrac{n}{N}}
\]
C'est un nombre entre $0$ et $1$, qu'on écrit souvent en pourcentage.
\end{definitionbox}

\begin{methodebox}[Méthode 6 --- Calculer une proportion en pourcentage]
\etapeexemples
\begin{questions}[label=\alph*)]
  \item Sur $20$ élèves, $9$ viennent au lycée à vélo. Quel pourcentage des élèves
        cela représente-t-il ?
  \item Un spectacle réunit $300$ spectateurs, dont $54$ enfants. Quelle est la
        proportion d'enfants, en pourcentage ?
\end{questions}
\begin{correction}
\cas{a)}
\begin{calculs}
p &= \dfrac{9}{20} &\qquad& \rem{la partie sur le total}\\
p &= \dfrac{45}{100} &\qquad& \rem{on multiplie en haut et en bas par $5$}\\
p &= \res{45\,\%} &&
\end{calculs}
\cas{b)}
\begin{calculs}
p &= \dfrac{54}{300} &\qquad& \rem{la partie sur le total}\\
p &= \dfrac{18}{100} &\qquad& \rem{on divise en haut et en bas par $3$}\\
p &= \res{18\,\%} &&
\end{calculs}
\end{correction}

\etape{À vous de jouer}
Écrire chaque proportion en pourcentage :
\[
  \dfrac{7}{20} \qquad\qquad \dfrac{1}{4} \qquad\qquad \dfrac{36}{200}
  \qquad\qquad \text{$18$ reçus sur $60$ candidats}
\]
\reponse{$35\,\%$ ;\; $25\,\%$ ;\; $18\,\%$ ;\; $30\,\%$.}
\end{methodebox}

\begin{definitionbox}[Définition 6 --- Taux d'évolution]
Une grandeur passe d'une valeur de départ $V_{0}$ à une valeur d'arrivée $V_{1}$. Son
\textbf{taux d'évolution} est
\[
  \boxed{t=\dfrac{V_{1}-V_{0}}{V_{0}}}
\]
\begin{itemize}
  \item $t>0$ : c'est une \textbf{hausse} ; $t<0$ : c'est une \textbf{baisse} ;
  \item on divise toujours par la valeur de \textbf{départ}.
\end{itemize}
\end{definitionbox}

\begin{methodebox}[Méthode 7 --- Calculer un taux d'évolution]
\etapeexemples
\begin{questions}[label=\alph*)]
  \item Un prix passe de $60$~€ à $69$~€. Calculer le taux d'évolution.
  \item Une récolte passe de $500$ à $470$ tonnes. Calculer le taux d'évolution.
  \item Une association passe de $600$ à $530$ membres. Justifier qu'elle a perdu plus
        de $10\,\%$ de ses membres.
\end{questions}
\begin{correction}
\cas{a)}
\begin{calculs}
t &= \dfrac{69-60}{60} &\qquad& \rem{arrivée moins départ, sur départ}\\
t &= \dfrac{9}{60} &&\\
t &= 0{,}15 &\qquad& \rem{$\dfrac{9}{60}=\dfrac{3}{20}=\dfrac{15}{100}$}
\end{calculs}
\explique{Le prix a \res{augmenté de $15\,\%$}.}{$t>0$ : hausse}
\cas{b)}
\begin{calculs}
t &= \dfrac{470-500}{500} &&\\
t &= \dfrac{-30}{500} &&\\
t &= -0{,}06 &\qquad& \rem{$\dfrac{30}{500}=\dfrac{6}{100}$}
\end{calculs}
\explique{La récolte a \res{baissé de $6\,\%$}.}{$t<0$ : baisse}
\cas{c)}
\explique{L'association a perdu $600-530=70$ membres.}{}
\explique{$10\,\%$ de $600$ font $60$, et $70>60$ : \res{elle a perdu plus de $10\,\%$ de
ses membres}.}{on compare la perte à $10\,\%$ du départ}
\explique{Autre rédaction : $600\times 0{,}9=540$ et $530<540$.}{comme dans les
corrigés de l'épreuve}
\end{correction}

\etape{À vous de jouer}
Calculer le taux d'évolution de chaque grandeur :
\begin{questions}
  \item un salaire qui passe de $2\,000$~€ à $2\,080$~€ ;
  \item une facture qui passe de $40$~€ à $30$~€ ;
  \item une production qui passe de $120$ à $150$ tonnes ;
  \item une ville passe de $4\,000$ à $4\,300$ habitants : a-t-elle gagné plus de
        $5\,\%$ d'habitants ?
\end{questions}
\reponse{$+4\,\%$ ;\; $-25\,\%$ ;\; $+25\,\%$ ;\; oui : $5\,\%$ de $4\,000$ font $200$,
et elle a gagné $300$ habitants.}
\end{methodebox}

\afaire{exercices 13 à 16}

\partiecours{Coefficient multiplicateur}

\begin{exemplebox}[Pour commencer]
Augmenter $150$ de $20\,\%$, c'est ajouter $0{,}2\times 150$ :
\[
  150+0{,}2\times 150=150\times(1+0{,}2) \qquad\text{soit}\qquad 150\times 1{,}2=180
\]
Le résultat s'obtient en \textbf{une seule multiplication}.
\end{exemplebox}

\begin{proprietebox}[Propriété 1 --- Coefficient multiplicateur]
\begin{itemize}
  \item Augmenter de $t\,\%$, c'est multiplier par $\boxed{1+\dfrac{t}{100}}$.
  \item Diminuer de $t\,\%$, c'est multiplier par $\boxed{1-\dfrac{t}{100}}$.
\end{itemize}
Ce nombre est le \textbf{coefficient multiplicateur} $C_{M}$ : $V_{1}=V_{0}\times C_{M}$.
\begin{itemize}
  \item $C_{M}>1$ : c'est une hausse ; $C_{M}<1$ : c'est une baisse.
  \item On lit le taux d'évolution sur $C_{M}$ : $t=C_{M}-1$.
\end{itemize}
\end{proprietebox}

\begin{center}
\renewcommand{\arraystretch}{1.4}
\begin{tabular}{|l|c|c|c|c|c|c|}
\hline
Évolution & $+30\,\%$ & $+4\,\%$ & $+100\,\%$ & $-20\,\%$ & $-6\,\%$ & $-50\,\%$\\
\hline
Coefficient multiplicateur & $1{,}3$ & $1{,}04$ & $2$ & $0{,}8$ & $0{,}94$ & $0{,}5$\\
\hline
\end{tabular}
\end{center}

\begin{methodebox}[Méthode 8 --- Appliquer une évolution avec le coefficient multiplicateur]
\etapeexemples
\begin{questions}[label=\alph*)]
  \item Un téléphone à $400$~€ baisse de $15\,\%$. Quel calcul donne son nouveau prix ?
  \item Une colonie de $2\,000$ abeilles augmente de $40\,\%$ au printemps. Combien
        d'abeilles compte-t-elle ensuite ?
  \item Un prix est multiplié par $0{,}94$, un autre par $1{,}15$. Quelles évolutions
        ont-ils subies ?
  \item Un billet passe de $80$~€ à $100$~€. Retrouver le taux d'évolution avec le
        coefficient multiplicateur.
\end{questions}
\begin{correction}
\cas{a)}
\explique{$C_{M}=1-0{,}15=0{,}85$}{baisse de $15\,\%$}
\explique{Nouveau prix : $400\times 0{,}85=\res{340}$~€.}{et non $400-0{,}15$}
\cas{b)}
\explique{$C_{M}=1+0{,}4=1{,}4$}{hausse de $40\,\%$}
\explique{$2\,000\times 1{,}4=\res{2\,800}$ abeilles.}{}
\cas{c)}
\explique{$0{,}94=1-0{,}06$ : \res{baisse de $6\,\%$}.}{$C_{M}<1$}
\explique{$1{,}15=1+0{,}15$ : \res{hausse de $15\,\%$}.}{$C_{M}>1$}
\cas{d)}
\begin{calculs}
C_{M} &= \dfrac{V_{1}}{V_{0}} &\qquad& \rem{arrivée sur départ}\\
C_{M} &= \dfrac{100}{80} &&\\
C_{M} &= 1{,}25 &\qquad& \rem{$1{,}25=1+0{,}25$}
\end{calculs}
\explique{Le prix a \res{augmenté de $25\,\%$}.}{même résultat qu'avec la méthode 7}
\end{correction}

\etape{À vous de jouer}
\begin{questions}
  \item Donner le coefficient multiplicateur de : $+8\,\%$ ; $-40\,\%$ ; $+200\,\%$.
  \item Donner l'évolution associée à : $\times 1{,}05$ ; $\times\dfrac{1}{2}$ ;
        $\times 2$.
  \item Un salaire de $1\,800$~€ augmente de $5\,\%$ ; une facture de $60$~€ baisse de
        $25\,\%$. Calculer les nouveaux montants.
  \item Un prix passe de $50$~€ à $40$~€. Calculer $C_{M}$, puis le taux d'évolution.
\end{questions}
\reponse{$1{,}08$ ;\; $0{,}6$ ;\; $3$ ;\; $+5\,\%$ ;\; $\frac{1}{2}=0{,}5$ : $-50\,\%$ ;\;
$+100\,\%$ ;\; $1\,890$~€ ;\; $45$~€ ;\; $C_{M}=0{,}8$, baisse de $20\,\%$.}
\end{methodebox}

\afaire{exercices 17 à 20}

\partiecours{Évolutions successives, évolution réciproque}

\begin{exemplebox}[Pour commencer]
Un jean à $80$~€ baisse de $20\,\%$, puis augmente de $20\,\%$ :
\[
  80\times 0{,}8=64 \qquad\text{puis}\qquad 64\times 1{,}2=76{,}80
\]
Il ne revient pas à $80$~€ : les pourcentages \textbf{ne s'ajoutent pas}.
\end{exemplebox}

\begin{proprietebox}[Propriété 2 --- Évolutions successives]
Le coefficient multiplicateur global de plusieurs évolutions successives est le
\textbf{produit} de leurs coefficients multiplicateurs :
\[
  \boxed{C_{M\,\text{global}}=C_{M1}\times C_{M2}\times\cdots}
\]
Une même évolution répétée $n$ fois a pour coefficient global $C_{M}^{n}$.
\end{proprietebox}

\begin{methodebox}[Méthode 9 --- Calculer un taux d'évolution global]
\etapeexemples
Calculer le taux d'évolution global dans chaque cas.
\begin{questions}[label=\alph*)]
  \item Un prix augmente de $20\,\%$, puis baisse de $20\,\%$.
  \item Le chiffre d'affaires d'un magasin augmente de $10\,\%$ en 2025, puis de
        $30\,\%$ en 2026.
  \item Un loyer augmente de $4\,\%$ par an pendant $3$ ans. On donne
        $1{,}04^{3}\approx 1{,}12$.
\end{questions}
\begin{correction}
\cas{a)}
\begin{calculs}
C_{M} &= 1{,}2\times 0{,}8 &\qquad& \rem{produit des coefficients}\\
C_{M} &= 0{,}96 &\qquad& \rem{$0{,}96=1-0{,}04$}
\end{calculs}
\explique{Globalement, le prix a \res{baissé de $4\,\%$}.}{}
\cas{b)}
\begin{calculs}
C_{M} &= 1{,}1\times 1{,}3 &&\\
C_{M} &= 1{,}43 &\qquad& \rem{$1{,}43=1+0{,}43$}
\end{calculs}
\explique{Le chiffre d'affaires a \res{augmenté de $43\,\%$}.}{et non de $40\,\%$}
\cas{c)}
\begin{calculs}
C_{M} &= 1{,}04^{3} &\qquad& \rem{la même évolution trois fois}\\
C_{M} &\approx 1{,}12 &&
\end{calculs}
\explique{Le loyer a \res{augmenté d'environ $12\,\%$} en $3$ ans.}{et non de $12\,\%$
exactement : $1{,}04^{3}=1{,}124\,864$}
\end{correction}

\etape{À vous de jouer}
Calculer le taux d'évolution global dans chaque cas :
\begin{questions}
  \item une hausse de $50\,\%$, puis une baisse de $40\,\%$ ;
  \item deux baisses successives de $10\,\%$ ;
  \item une hausse de $20\,\%$ par an pendant $2$ ans.
\end{questions}
\reponse{$1{,}5\times 0{,}6=0{,}9$ : $-10\,\%$ ;\; $0{,}9\times 0{,}9=0{,}81$ :
$-19\,\%$ ;\; $1{,}2^{2}=1{,}44$ : $+44\,\%$.}
\end{methodebox}

\begin{definitionbox}[Définition 7 --- Évolution réciproque]
L'\textbf{évolution réciproque} d'une évolution est celle qui ramène à la valeur de
départ. Deux évolutions successives se \textbf{compensent} quand le produit de leurs
coefficients multiplicateurs vaut $1$ :
\[
  \boxed{C_{M1}\times C_{M2}=1}
\]
\end{definitionbox}

\begin{remarquebox}[Cas à connaître]
\begin{center}
\renewcommand{\arraystretch}{1.4}
\begin{tabular}{|l|c|c|c|}
\hline
Évolution & $+25\,\%$ & $+100\,\%$ (doubler) & $+300\,\%$ (quadrupler)\\
\hline
Compensée par & $-20\,\%$ & $-50\,\%$ (moitié) & $-75\,\%$ (quart)\\
\hline
Vérification & $1{,}25\times 0{,}8=1$ & $2\times 0{,}5=1$ & $4\times 0{,}25=1$\\
\hline
\end{tabular}
\end{center}
Une hausse de $t\,\%$ n'est \textbf{jamais} compensée par une baisse de $t\,\%$.
\end{remarquebox}

\begin{methodebox}[Méthode 10 --- Vérifier qu'une évolution en compense une autre]
\etapeexemples
\begin{questions}[label=\alph*)]
  \item Un prix double. Une baisse de $50\,\%$ le ramène-t-elle à sa valeur de départ ?
  \item Une hausse de $20\,\%$ est-elle compensée par une baisse de $20\,\%$ ?
  \item Une action perd $20\,\%$ de sa valeur. Pour la retrouver, il faut une hausse
        de : $20\,\%$ ; $25\,\%$ ; $40\,\%$ ?
\end{questions}
\begin{correction}
\cas{a)}
\begin{calculs}
C_{M1}\times C_{M2} &= 2\times 0{,}5 &\qquad& \rem{doubler : $\times 2$}\\
C_{M1}\times C_{M2} &= 1 &&
\end{calculs}
\explique{\res{Oui} : la baisse de $50\,\%$ compense le doublement.}{le produit vaut $1$}
\cas{b)}
\begin{calculs}
C_{M1}\times C_{M2} &= 1{,}2\times 0{,}8 &&\\
C_{M1}\times C_{M2} &= 0{,}96 &\qquad& \rem{différent de $1$}
\end{calculs}
\explique{\res{Non} : on n'est pas revenu à la valeur de départ.}{il manque $4\,\%$}
\cas{c)}
\explique{$0{,}8\times 1{,}2=0{,}96$ ; \ $0{,}8\times 1{,}25=1$ ; \
$0{,}8\times 1{,}4=1{,}12$.}{on teste chaque réponse}
\explique{Il faut une \res{hausse de $25\,\%$}.}{seul produit égal à $1$}
\end{correction}

\etape{À vous de jouer}
\begin{questions}
  \item Une baisse de $50\,\%$ est-elle compensée par une hausse de $100\,\%$ ?
  \item Une hausse de $10\,\%$ est-elle compensée par une baisse de $10\,\%$ ?
  \item Un prix est multiplié par $4$. Pour revenir au prix de départ, il faut une
        baisse de : $75\,\%$ ; $300\,\%$ ; $25\,\%$ ?
\end{questions}
\reponse{oui, $0{,}5\times 2=1$ ;\; non, $1{,}1\times 0{,}9=0{,}99$ ;\; $75\,\%$, car
$4\times 0{,}25=1$.}
\end{methodebox}

\afaire{exercices 21 à 25}

\partiecours{Fonctions affines : image, antécédent}

\begin{definitionbox}[Définition 8 --- Fonction affine]
Une fonction $f$ est \textbf{affine} quand il existe deux nombres $a$ et $b$ tels que,
pour tout nombre réel $x$ :
\[
  \boxed{f(x)=ax+b}
\]
$a$ est le \textbf{coefficient directeur}, $b$ l'\textbf{ordonnée à l'origine}.
\begin{itemize}
  \item Si $b=0$, $f(x)=ax$ : $f$ est \textbf{linéaire} (proportionnalité).
  \item Si $a=0$, $f(x)=b$ : $f$ est \textbf{constante}.
\end{itemize}
\end{definitionbox}

\begin{exemplebox}[Exemples]
\begin{itemize}
  \item Une course en VTC coûte $5$~€ de prise en charge, puis $2$~€ par kilomètre.
        Pour $x$~km, on paie $f(x)=2x+5$ : $a=2$ et $b=5$.
  \item Baisser un prix $x$ de $10\,\%$ donne $h(x)=0{,}9x$ : $h$ est linéaire.
\end{itemize}
\end{exemplebox}

\begin{remarquebox}[Antécédent par une fonction affine]
Pour une fonction affine non constante, chaque nombre a un seul antécédent.
\end{remarquebox}

\begin{methodebox}[Méthode 11 --- Calculer une image]
\etapeexemples
\begin{questions}[label=\alph*)]
  \item Soit $f(x)=-3x+8$. Calculer $f(2)$ et $f(-5)$.
  \item VTC : calculer le prix d'une course de $12$~km.
\end{questions}
\begin{correction}
\cas{a)}
\begin{calculs}
f(2) &= -3\times 2+8 &\qquad& \rem{on remplace $x$ par $2$}\\
f(2) &= \res{2} &&\\
f(-5) &= -3\times(-5)+8 &\qquad& \rem{parenthèses autour de $-5$}\\
f(-5) &= \res{23} &\qquad& \rem{$15+8$}
\end{calculs}
\cas{b)}
\begin{calculs}
f(12) &= 2\times 12+5 &&\\
f(12) &= 29 &&
\end{calculs}
\explique{La course coûte $\res{29}$~€.}{on répond dans le contexte}
\end{correction}

\etape{À vous de jouer}
\begin{questions}
  \item Soit $f(x)=4x-7$. Calculer $f(0)$, $f(3)$ et $f(-2)$.
  \item Soit $h(x)=0{,}5x+3$. Calculer $h(10)$ et $h(-4)$.
  \item VTC : calculer le prix d'une course de $8$~km.
\end{questions}
\reponse{$-7$ ;\; $5$ ;\; $-15$ ;\; $8$ ;\; $1$ ;\; $21$~€.}
\end{methodebox}

\begin{methodebox}[Méthode 12 --- Calculer un antécédent]
\etapeexemples
\begin{questions}[label=\alph*)]
  \item Soit $f(x)=4x-7$. Calculer l'antécédent de $9$.
  \item VTC : une course a coûté $35$~€. Quelle distance a-t-on parcourue ?
\end{questions}
\begin{correction}
\cas{a)}
\begin{calculs}
4x-7 &= 9 &\qquad& \rem{on écrit $f(x)=9$}\\
4x &= 16 &\qquad& \rem{on ajoute $7$}\\
x &= 4 &\qquad& \rem{on divise par $4$}
\end{calculs}
\explique{L'antécédent de $9$ est $\res{4}$. Contrôle : $f(4)=16-7=9$.}{}
\cas{b)}
\begin{calculs}
2x+5 &= 35 &\qquad& \rem{on écrit $f(x)=35$}\\
2x &= 30 &\qquad& \rem{on retire $5$}\\
x &= 15 &\qquad& \rem{on divise par $2$}
\end{calculs}
\explique{On a parcouru $\res{15}$~km.}{}
\end{correction}

\etape{À vous de jouer}
\begin{questions}
  \item Soit $f(x)=-3x+8$. Calculer l'antécédent de $14$, puis celui de $2$.
  \item VTC : une course a coûté $45$~€. Quelle distance a-t-on parcourue ?
  \item Soit $h(x)=5x+2$. Résoudre $h(x)=-8$.
\end{questions}
\reponse{$-2$ ;\; $2$ ;\; $2x=40$ donc $20$~km ;\; $x=-2$.}
\end{methodebox}

\afaire{exercices 26 à 29}

\partiecours{Trouver une fonction affine}

\begin{proprietebox}[Propriété 3 --- Coefficient directeur à partir de deux points]
Si la droite qui représente $f(x)=ax+b$ passe par $A(x_{A}\,;y_{A})$ et
$B(x_{B}\,;y_{B})$, alors :
\[
  \boxed{a=\dfrac{y_{B}-y_{A}}{x_{B}-x_{A}}}
  \qquad\text{(différence des ordonnées sur différence des abscisses)}
\]
Quand $x$ augmente de $1$, $f(x)$ augmente de $a$ (diminue si $a<0$) ; et $b=f(0)$ est
la valeur de départ.
\end{proprietebox}

\begin{methodebox}[Méthode 13 --- Déterminer une fonction affine]
\etapeexemples
\begin{questions}[label=\alph*)]
  \item Un abonnement de vidéo coûte $9$~€ par mois, plus $3$~€ par film loué.
        Exprimer le prix mensuel $p(x)$ pour $x$ films loués.
  \item La droite qui représente $f$ passe par $A(1\,;4)$ et $B(4\,;13)$. Déterminer
        $f(x)$.
  \item La droite qui représente $g$ passe par $A(-2\,;5)$ et $B(-4\,;9)$. Déterminer
        $g(x)$.
  \item Une facture d'électricité vaut $70$~€ pour $200$~kWh et $100$~€ pour
        $300$~kWh. On la modélise par une fonction affine $F$ de la consommation $x$.
        Déterminer $F(x)$, puis interpréter $a$ et $b$.
\end{questions}
\begin{correction}
\cas{a)}
\explique{$b=9$ et $a=3$ : $p(x)=\res{3x+9}$.}{$b$ : la partie fixe ; $a$ : le prix d'un
film}
\cas{b)}
\begin{calculs}
a &= \dfrac{y_{B}-y_{A}}{x_{B}-x_{A}} &&\\
a &= \dfrac{13-4}{4-1} &&\\
a &= 3 &\qquad& \rem{$\dfrac{9}{3}$}\\
3\times 1+b &= 4 &\qquad& \rem{le point $A$ est sur la droite}\\
b &= 1 &\qquad& \rem{on retire $3$}
\end{calculs}
\explique{$f(x)=\res{3x+1}$. Contrôle avec $B$ : $f(4)=12+1=13$.}{}
\cas{c)}
\begin{calculs}
a &= \dfrac{9-5}{-4-(-2)} &\qquad& \rem{$-4-(-2)=-4+2=-2$}\\
a &= \dfrac{4}{-2} &&\\
a &= -2 &&\\
-2\times(-2)+b &= 5 &\qquad& \rem{avec le point $A$}\\
4+b &= 5 &&\\
b &= 1 &&
\end{calculs}
\explique{$g(x)=\res{-2x+1}$. Contrôle avec $B$ : $g(-4)=8+1=9$.}{}
\cas{d)}
\begin{calculs}
a &= \dfrac{100-70}{300-200} &\qquad& \rem{points $(200\,;70)$ et $(300\,;100)$}\\
a &= 0{,}3 &\qquad& \rem{$\dfrac{30}{100}$}\\
0{,}3\times 200+b &= 70 &&\\
b &= 10 &\qquad& \rem{$60+b=70$}
\end{calculs}
\explique{$F(x)=\res{0{,}3x+10}$ : \res{le kWh coûte $0{,}30$~€ et l'abonnement $10$~€}.}{}
\end{correction}

\etape{À vous de jouer}
\begin{questions}
  \item Un club de danse coûte $40$~€ d'inscription, puis $20$~€ par mois. Exprimer le
        coût $C(x)$ pour $x$ mois.
  \item La droite qui représente $g$ passe par $A(2\,;9)$ et $B(5\,;3)$. Déterminer
        $g(x)$.
  \item La droite qui représente $h$ passe par $A(-1\,;-4)$ et $B(3\,;8)$. Déterminer
        $h(x)$.
\end{questions}
\reponse{$C(x)=20x+40$ ;\; $a=\dfrac{3-9}{5-2}=-2$ et $g(x)=-2x+13$ ;\;
$a=\dfrac{8-(-4)}{3-(-1)}=3$ et $h(x)=3x-1$.}
\end{methodebox}

\afaire{exercices 30 à 32}

\partiecours{Tracer une fonction affine}

\begin{proprietebox}[Propriété 4 --- Représentation graphique et variations]
La représentation graphique de $f(x)=ax+b$ est une \textbf{droite}, d'\textbf{équation
réduite} $\boxed{y=ax+b}$ : le point $M(x\,;y)$ est sur la droite quand $y=ax+b$.
\begin{itemize}
  \item elle coupe l'axe des ordonnées au point $(0\,;b)$ ;
  \item pour un déplacement horizontal de $1$, le déplacement vertical vaut $a$ (on
        monte si $a>0$, on descend si $a<0$).
\end{itemize}
Sens de variation : $a>0$, $f$ est \textbf{croissante} ; $a<0$, $f$ est
\textbf{décroissante} ; $a=0$, $f$ est \textbf{constante}.
\end{proprietebox}

\begin{methodebox}[Méthode 14 --- Tracer la droite d'une fonction affine]
\etapeexemples
Tracer les droites qui représentent $f(x)=3x-2$ et $g(x)=-0{,}5x+3$.
\begin{correction}
\cas{Résolution}\par\nopagebreak
\noindent
\begin{minipage}[c]{0.58\linewidth}\raggedright
\explique{$f(0)=-2$ et $f(2)=4$ : points $(0\,;-2)$ et $(2\,;4)$.}{deux images
sur le quadrillage}
\explique{$g(0)=3$ et $g(4)=1$ : points $(0\,;3)$ et $(4\,;1)$.}{abscisse paire : pas
de demi-unité}
\explique{On place les croix et on trace les droites.}{}
\explique{$a=3>0$ : la droite de $f$ monte ; $a=-0{,}5<0$ : celle de $g$
descend.}{contrôle}
\end{minipage}\hfill
\begin{minipage}[c]{0.38\linewidth}\centering
\begin{tikzpicture}
\begin{axis}[repere,
    x=0.42cm,y=0.42cm,
    xmin=-2.5,xmax=5.5,ymin=-3.5,ymax=5.5,
    xtick={-2,...,5},ytick={-3,...,5},
    xlabel={$x$},ylabel={$y$}]
  \addplot[coulCourbe,line width=1.1pt,domain=-0.45:2.45]{3*x-2};
  \addplot[coulCourbeB,line width=1.1pt,domain=-2.5:5.5]{-0.5*x+3};
  \addplot[only marks,mark=x,mark options={line width=0.6pt},mark size=2pt]
    coordinates {(0,-2)(2,4)(0,3)(4,1)};
  \node[coulCourbe,right] at (axis cs:2.3,4.6) {$f$};
  \node[coulCourbeB,above] at (axis cs:-1.8,4) {$g$};
\end{axis}
\end{tikzpicture}
\end{minipage}
\end{correction}

\etape{À vous de jouer}
Tracer dans un repère les droites qui représentent $h(x)=x-2$ ; $k(x)=-2x+3$ et
$m(x)=\dfrac{1}{2}x+1$.
\reponse{$h$ : $(0\,;-2)$ et $(3\,;1)$ ;\; $k$ : $(0\,;3)$ et $(2\,;-1)$ ;\; $m$ :
$(0\,;1)$ et $(2\,;2)$.}
\end{methodebox}

\begin{proprietebox}[Propriété 5 --- Fonctions linéaires et proportionnalité]
Une fonction \textbf{linéaire} $f(x)=ax$ (cas $b=0$) traduit une situation de
\textbf{proportionnalité} : $a$ est le coefficient de proportionnalité, la quantité
pour une unité de $x$.
\begin{itemize}
  \item Sa droite passe par l'\textbf{origine} du repère.
  \item Appliquer une évolution de $t\,\%$ est une fonction linéaire :
        $x\mapsto C_{M}\times x$.
\end{itemize}
\end{proprietebox}

\begin{methodebox}[Méthode 15 --- Utiliser une fonction linéaire (proportionnalité)]
\etapeexemples
\begin{questions}[label=\alph*)]
  \item Des céréales contiennent $8$~g de fibres pour $100$~g. Quelle masse de fibres
        contient un paquet de $350$~g ?
  \item Un mile vaut environ $1{,}6$~km. Convertir $50$ miles en kilomètres, puis
        $24$~km en miles.
  \item Parmi $f(x)=-3x$ ; $g(x)=3x-1$ et $h(x)=0{,}4x$, lesquelles sont linéaires ?
\end{questions}
\begin{correction}
\cas{a)}
\explique{$a=\dfrac{8}{100}=0{,}08$ : dans $x$ grammes, il y a $f(x)=0{,}08x$ grammes
de fibres.}{la quantité pour $1$~g}
\explique{$f(350)=0{,}08\times 350=28$ : le paquet contient $\res{28}$~g de fibres.}{}
\cas{b)}
\explique{$f(x)=1{,}6x$ donne les kilomètres pour $x$ miles : $f(50)=\res{80}$~km.}{}
\begin{calculs}
1{,}6x &= 24 &\qquad& \rem{on cherche l'antécédent de $24$}\\
x &= 15 &\qquad& \rem{$\dfrac{24}{1{,}6}=\dfrac{240}{16}$}
\end{calculs}
\explique{$24$~km font environ $\res{15}$ miles.}{}
\cas{c)}
\explique{\res{$f$ et $h$ sont linéaires} : leurs droites passent par l'origine.}{$b=0$}
\explique{$g$ est affine, mais pas linéaire.}{$b=-1$}
\end{correction}

\etape{À vous de jouer}
\begin{questions}
  \item Une imprimante imprime $20$ pages par minute. Exprimer le nombre $P(t)$ de pages
        imprimées en $t$ minutes ; calculer $P(4)$, puis le temps pour imprimer
        $300$ pages.
  \item En soldes, un prix $x$ baisse de $15\,\%$. Exprimer le prix soldé $R(x)$ ;
        calculer $R(60)$, puis le prix de départ d'un article soldé à $34$~€.
  \item Tracer la droite qui représente $k(x)=3x$.
\end{questions}
\reponse{$P(t)=20t$ ; $P(4)=80$ ; $20t=300$ donc $15$~min ;\; $R(x)=0{,}85x$ ;
$R(60)=51$~€ ; $0{,}85x=34$ donc $40$~€ ;\; droite passant par $(0\,;0)$ et $(1\,;3)$.}
\end{methodebox}

\afaire{exercices 33 à 36}

\partiecours{Lire l'équation d'une droite}

\begin{proprietebox}[Propriété 6 --- Lire $a$ et $b$ sur la droite]
Sur la droite d'équation réduite $y=ax+b$ :
\begin{itemize}
  \item $b$ est l'ordonnée du point où la droite coupe l'axe des ordonnées ;
  \item entre deux points de la droite,
        $\boxed{a=\dfrac{\text{déplacement vertical}}{\text{déplacement horizontal}}}$,
        le déplacement vertical étant négatif quand on descend (c'est
        $\dfrac{y_{B}-y_{A}}{x_{B}-x_{A}}$ lu sur le quadrillage).
\end{itemize}
\end{proprietebox}

\begin{methodebox}[Méthode 16 --- Lire l'équation réduite d'une droite]
\etapeexemples
\noindent
\begin{minipage}[c]{0.58\linewidth}\raggedright
Les droites $(d_{1})$ et $(d_{2})$ représentent deux fonctions affines $u$ et $v$.
Donner l'équation réduite de chaque droite, puis le sens de variation de $u$ et de $v$.
\end{minipage}\hfill
\begin{minipage}[c]{0.38\linewidth}\centering
\begin{tikzpicture}
\begin{axis}[repere,
    x=0.42cm,y=0.42cm,
    xmin=-1.5,xmax=6.5,ymin=-3.5,ymax=4.5,
    xtick={-1,...,6},ytick={-3,...,4},
    xlabel={$x$},ylabel={$y$}]
  \addplot[coulCourbe,line width=1.1pt,domain=-0.6:3.2]{-2*x+3};
  \addplot[coulCourbeB,line width=1.1pt,domain=-1.5:6.5]{x/3-1};
  \node[coulCourbe,right] at (axis cs:-0.5,4) {$(d_{1})$};
  \node[coulCourbeB,above] at (axis cs:5.4,1) {$(d_{2})$};
\end{axis}
\end{tikzpicture}
\end{minipage}

\begin{correction}
\cas{$(d_{1})$}
\explique{$b=3$}{la droite coupe l'axe des ordonnées en $(0\,;3)$}
\explique{De $(0\,;3)$ à $(1\,;1)$ : on avance de $1$ et on descend de $2$, donc
$a=\dfrac{-2}{1}=-2$.}{vertical sur horizontal}
\explique{$(d_{1})$ : $y=\res{-2x+3}$, et $u$ est \res{décroissante}.}{$a<0$}
\cas{$(d_{2})$}
\explique{$b=-1$}{point $(0\,;-1)$}
\explique{De $(0\,;-1)$ à $(3\,;0)$ : on avance de $3$ et on monte de $1$, donc
$a=\dfrac{1}{3}$.}{}
\explique{$(d_{2})$ : $y=\res{\dfrac{1}{3}x-1}$, et $v$ est \res{croissante}.}{$a>0$}
\end{correction}

\etape{À vous de jouer}
Une droite passe par $(0\,;-3)$ et $(1\,;1)$ ; une autre par $(0\,;2)$ et $(2\,;1)$.
Donner l'équation réduite de chaque droite et le sens de variation de la fonction
qu'elle représente.
\reponse{$y=4x-3$, croissante ;\; $y=-0{,}5x+2$, décroissante.}
\end{methodebox}

\afaire{exercices 37 à 39}

\partiecours{Comparer deux fonctions affines}

\begin{remarquebox}[Rappel --- Résoudre une équation $ax+b=cx+d$]
On regroupe les termes en $x$ d'un côté et les nombres de l'autre, puis on divise par le
coefficient de $x$ :
\begin{calculs}
7x+2 &= 3x+14 &&\\
4x &= 12 &\qquad& \rem{on retire $3x$ et $2$}\\
x &= 3 &\qquad& \rem{on divise par $4$}
\end{calculs}
\end{remarquebox}

\begin{methodebox}[Méthode 17 --- Comparer deux offres]
\etapeexemples
\noindent
\begin{minipage}[c]{0.58\linewidth}\raggedright
Une agence loue une camionnette à la journée. Tarif A : $20$~€, plus $0{,}50$~€ par
km ; tarif B : $80$~€, plus $0{,}20$~€ par km. On note $f(x)$ et $g(x)$ les prix pour
$x$~km. Quel tarif choisir selon la distance ? Le graphique ci-contre représente les
deux tarifs.
\end{minipage}\hfill
\begin{minipage}[c]{0.38\linewidth}\centering
\begin{tikzpicture}
\begin{axis}[repere,
    x=0.016cm,y=0.016cm,
    xmin=0,xmax=330,ymin=0,ymax=200,
    xtick={50,100,...,300},ytick={20,40,...,180},
    xticklabels={,100,,200,,300},yticklabels={,40,,80,,120,,160,},
    xlabel={km},ylabel={€}]
  \addplot[coulCourbe,line width=1.1pt,domain=0:320]{0.5*x+20};
  \addplot[coulCourbeB,line width=1.1pt,domain=0:320]{0.2*x+80};
  \addplot[only marks,mark=x,mark options={line width=0.6pt},mark size=2pt]
    coordinates {(200,120)};
  \node[coulCourbe,left] at (axis cs:280,170) {A};
  \node[coulCourbeB,below] at (axis cs:300,140) {B};
\end{axis}
\end{tikzpicture}
\end{minipage}

\begin{correction}
\cas{Résolution}\par\nopagebreak
\explique{$f(x)=0{,}5x+20$ et $g(x)=0{,}2x+80$.}{partie fixe, plus le prix d'un km}
\explique{Les droites se coupent au point $(200\,;120)$ ; avant $200$~km, la droite de
A est en dessous.}{lecture graphique}
\begin{calculs}
0{,}5x+20 &= 0{,}2x+80 &\qquad& \rem{même prix pour les deux tarifs}\\
0{,}3x &= 60 &\qquad& \rem{on retire $0{,}2x$ et $20$}\\
x &= 200 &\qquad& \rem{$\dfrac{60}{0{,}3}=\dfrac{600}{3}$}
\end{calculs}
\explique{\res{Pour moins de $200$~km, le tarif A est le moins cher ; au-delà, le tarif B}.
Pour $200$~km, les deux coûtent $120$~€.}{on conclut dans le contexte}
\end{correction}

\etape{À vous de jouer}
\begin{questions}
  \item Résoudre $4x+3=x+12$ ; puis $6x-5=2x+11$ ; puis $2x+5=-4x+8$.
  \item Salle A : $20$~€ d'inscription et $30$~€ par mois ; salle B : $100$~€
        d'inscription et $20$~€ par mois. À partir de combien de mois la salle B
        est-elle la moins chère ?
\end{questions}
\reponse{$x=3$ ;\; $x=4$ ;\; $6x=3$ donc $x=0{,}5$ ;\; $30x+20=20x+100$ donc $x=8$ :
au-delà de $8$ mois.}
\end{methodebox}

\afaire{exercices 40 à 42}

\partiecours{Fonction carré}

\begin{definitionbox}[Définition 9 --- Fonction carré]
La \textbf{fonction carré} est définie sur $\R$ par $f(x)=x^{2}$. Sa courbe est une
\textbf{parabole} de sommet l'origine $O$.
\end{definitionbox}

\begin{center}
\begin{minipage}{0.40\linewidth}\centering
\begin{tikzpicture}[x=0.42cm,y=0.42cm]
  \repere{-3}{-1}{3}{9}
  \gradx{-2,-1,1,2} \grady{1,2,3,4,5,6,7,8}
  \draw[coulCourbe,line width=1.1pt,domain=-3:3,samples=60] plot (\x,{\x*\x});
  \foreach \p in {(-3,9),(-2,4),(-1,1),(0,0),(1,1),(2,4),(3,9)} \path[coulCourbe] \p node[croix]{};
\end{tikzpicture}
\end{minipage}\hfill
\begin{minipage}{0.55\linewidth}\centering
\begin{tikzpicture}[scale=0.85]
  \tkzTabInit[lgt=1.6,espcl=1.8]{$x$/1,$x^{2}$/2}{$-\infty$,$0$,$+\infty$}
  \tkzTabVar{+/$+\infty$,-/$0$,+/$+\infty$}
\end{tikzpicture}

\begin{proprietebox}[Propriété 7 --- Carré]
Décroissante sur $\left]-\infty\,;0\right]$, croissante sur $\left[0\,;+\infty\right[$.
Un carré est toujours positif ou nul, et deux nombres opposés ont le même carré :
$(-3)^{2}=3^{2}=9$.
\end{proprietebox}
\end{minipage}
\end{center}

\begin{proprietebox}[Propriété 8 --- Équation $x^{2}=k$]
Si $k>0$ : deux solutions, $-\sqrt{k}$ et $\sqrt{k}$. Si $k=0$ : une solution, $0$. Si
$k<0$ : aucune solution.
\end{proprietebox}

\begin{methodebox}[Méthode 18 --- Résoudre une équation en isolant $x^{2}$]
\etapeexemples
Résoudre dans $\R$ les équations suivantes.
\begin{questions}[label=\alph*)]
  \item $x^{2}=16$
  \item $5x^{2}+1=2x^{2}+16$
  \item $3x^{2}+8=5$
\end{questions}
\begin{correction}
\cas{a)}
\explique{$S=\res{\ens{-4\,;4}}$.}{$16>0$ : deux solutions ; on n'oublie pas $-4$, car
$(-4)^{2}=16$}
\cas{b)}
\begin{calculs}
5x^{2}+1 &= 2x^{2}+16 &&\\
5x^{2}-2x^{2} &= 16-1 &\qquad& \rem{on isole $x^{2}$ comme on isole $x$}\\
3x^{2} &= 15 &&\\
x^{2} &= 5 &&
\end{calculs}
\explique{$S=\res{\ens{-\sqrt{5}\,;\sqrt{5}}}$.}{$5>0$ : deux solutions}
\cas{c)}
\begin{calculs}
3x^{2}+8 &= 5 &&\\
3x^{2} &= -3 &&\\
x^{2} &= -1 &&
\end{calculs}
\explique{$S=\res{\varnothing}$.}{un carré n'est jamais négatif}
\end{correction}

\etape{À vous de jouer}
Résoudre dans $\R$ :
\[
  x^{2}=64 \qquad 4x^{2}-5=2x^{2}+9 \qquad 3x^{2}+1=13 \qquad 2x^{2}+9=1
\]
\reponse{$S=\{-8\,;8\}$ ;\; $x^{2}=7$, $S=\left\{-\sqrt{7}\,;\sqrt{7}\right\}$ ;\;
$x^{2}=4$, $S=\{-2\,;2\}$ ;\; $x^{2}=-4$, $S=\varnothing$.}
\end{methodebox}

\afaire{exercices 43 à 45, puis les exercices de type épreuve 46 à 48}

\end{document}
